% =====================================================================
%  Alinhamento múltiplo de sequências
%  Prática 4 -- Disciplina Métodos Computacionais (UNESP)
%
%  Compilar com pdfLaTeX (Overleaf: Menu > Compiler > pdfLaTeX).
%
%  VERSÃO: por padrão gera a versão do ALUNO (sem respostas).
%  Para a versão do PROFESSOR (com gabarito), troque \professorfalse
%  por \professortrue na linha indicada abaixo, ou compile com
%     pdflatex "\def\versaoprofessor{}\input{apostila_alinhamento.tex}"
% =====================================================================
\documentclass[11pt,a4paper]{article}

\newif\ifprofessor
\professorfalse            % <-- troque para \professortrue para incluir o gabarito
\ifdefined\versaoprofessor \professortrue \fi

\usepackage[T1]{fontenc}
\usepackage[utf8]{inputenc}
\usepackage{lmodern}        % Latin Modern: versão vetorial da Computer Modern
\usepackage{textcomp}
\usepackage[provide=*,brazilian]{babel}
\usepackage[a4paper,margin=2.5cm]{geometry}
\usepackage{fancyvrb}
\usepackage{upquote}
\usepackage{array}
\usepackage{booktabs}
\usepackage{longtable}
\usepackage{enumitem}
\usepackage{comment}
\usepackage{url}
\usepackage[hidelinks]{hyperref}

\fvset{xleftmargin=1.5em,fontsize=\small}
\setlist{itemsep=0.2em,topsep=0.4em}
\setlength{\parskip}{0.35em}
\setlength{\emergencystretch}{3em}
\urlstyle{tt}

\newcommand{\arq}[1]{\path{#1}}
\newcommand{\ui}[1]{\textit{#1}}
\newcommand{\tecla}[1]{\textsf{#1}}

\ifprofessor
  \specialcomment{resposta}{\begin{quote}\small\textbf{Gabarito.}\ }{\end{quote}}
\else
  \excludecomment{resposta}
\fi

\title{Alinhamento múltiplo de sequências\\[0.3em]
\large Prática 4: MUSCLE no MEGA, MAFFT no terminal e concatenação no R}
\author{Disciplina Métodos Computacionais\\Universidade Estadual Paulista (UNESP)}
\date{\ifprofessor Versão do professor (com gabarito)\else Versão do aluno\fi\\ Novembro de 2026}

\begin{document}
\maketitle

\paragraph{Origem do material.} Esta apostila é uma versão revisada do tutorial
\emph{Alinhamento}, publicado em \url{https://pedrotaucce.github.io/filogenia/}. O alinhamento
com o MAFFT passou a ser feito no terminal, os dois marcadores são alinhados de uma só vez com
um laço, e a prática termina com a concatenação dos dois alinhamentos no R. As mudanças estão
listadas no Apêndice~\ref{ap:mudancas}.

\paragraph{Arquivos.} A pasta \arq{pratica_alinhamento} contém:
\begin{itemize}
\item \arq{sequencias/16S_ids.fasta} e \arq{sequencias/COI_ids.fasta}: as matrizes baixadas
por número de acesso na Prática~3 (22 sequências de 16S e 14 de COI), com nomes no formato
\verb|Genero_especie_acesso|. Se preferir, substitua-as pelos arquivos que você mesmo baixou;
\item \arq{R/04_concatenar.R}: o script da Seção~\ref{sec:concatenar}.
\end{itemize}

\tableofcontents

% =====================================================================
\section{Alinhamento e homologia primária}
% =====================================================================

Apesar de muitos sistematas darem pouca importância ao alinhamento e o tratarem apenas como
parte do processo para a obtenção de uma árvore filogenética, esta etapa é extremamente
importante para a análise filogenética como um todo, pois é nela que definimos a
\textbf{hipótese de homologia primária} entre as nossas sequências. Em outras palavras, nesta
etapa assumimos que as bases alinhadas em uma mesma coluna \textbf{divergiram de um mesmo
estado ancestral}.

Sendo esta uma etapa tão importante, existem vários programas e algoritmos comumente
utilizados para alinhar sequências de DNA (e também de aminoácidos, que não abordaremos no
curso), como o ClustalW \cite{larkin2007}, o T-Coffee \cite{notredame2000}, o MUSCLE
\cite{edgar2004} e o MAFFT \cite{katoh2013}. Nesta prática usaremos os dois últimos.

Continuamos o trabalho com a filogenia de Hylodidae da Prática~3. Agora que você adquiriu as
sequências de 16S rRNA e COI no GenBank, precisa alinhá-las. Usaremos as matrizes baixadas
por número de acesso, menores e mais controladas, com um exemplar por espécie. Elas incluem
\emph{Limnomedusa macroglossa} (Alsodidae), que servirá como raiz da árvore nas próximas
análises.

% =====================================================================
\section{MUSCLE no MEGA 12}\label{sec:mega}
% =====================================================================

O MEGA \cite{kumar2024} traz o MUSCLE embutido. Vamos usá-lo para alinhar o COI:
\begin{enumerate}
\item No MEGA 12, vá em \ui{Align > Edit/Build Alignment}, escolha \ui{Retrieve sequences
from a file} e abra \arq{sequencias/COI_ids.fasta}. Está tudo desalinhado.
\item Selecione todas as sequências (\ui{Edit > Select All}, ou \tecla{Ctrl+A}) e vá em
\ui{Alignment > Align by MUSCLE}. Se o MEGA oferecer as opções \ui{DNA} e \ui{Codons},
escolha \ui{DNA}.
\item Na janela de parâmetros do MUSCLE, mantenha os valores padrão e clique em \ui{Compute}.
Aguarde um pouquinho e pronto: as sequências estão alinhadas.
\end{enumerate}

O COI é um gene codificante, então é interessante verificar se as sequências estão corretas:
se começam no códon certo e se os possíveis \emph{gaps} são múltiplos de três, por exemplo.
Isso influenciará etapas posteriores, por isso é interessante que o alinhamento comece sempre
na posição~1 do códon. A tradução também serve para checar se há \emph{stop codons} no meio
das sequências, o que poderia indicar que você sequenciou um pseudogene ou o gene errado.

\begin{enumerate}[resume]
\item Clique na aba \ui{Translated Protein Sequences}. Quando aparecer a pergunta
\ui{The Current Genetic Code is: Standard. Is this correct?}, clique em \ui{No}.
\item Na janela seguinte, como se trata de um fragmento mitocondrial de anfíbios, marque
\ui{Vertebrate Mitochondrial} e confirme.
\item Verifique se há \emph{stop codons} (indicados pelo símbolo \texttt{*}) em alguma
posição. Se houver, provavelmente o alinhamento não começa na posição~1 do códon, ou o gene
sequenciado não é o que você queria.
\end{enumerate}

\begin{resposta}
Alinhamento com cerca de 712 colunas (712 com o MUSCLE 3.8.31), sem \emph{gaps} internos:
as sequências só diferem no início e no fim. Traduzido com o código mitocondrial de
vertebrados a partir da coluna~1, \textbf{não há nenhum \texttt{*}}. Se os alunos olharem a
aba de proteínas \emph{antes} de alinhar, verão \texttt{*} em 10 das 14 sequências, porque
cada uma começa em uma posição diferente do códon; só \emph{C.~caramaschii}, \emph{H.~ornatus},
\emph{H.~sazimai} e \emph{P.~boticariana} começam na posição~1. Com o código \ui{Standard},
aparecem \texttt{*} mesmo no alinhamento, porque \texttt{TGA} codifica triptofano no código
mitocondrial de vertebrados.
\end{resposta}

% =====================================================================
\section{MAFFT no terminal}\label{sec:mafft}
% =====================================================================

O MAFFT é um programa um pouco mais completo que o MUSCLE, por ter vários algoritmos
diferentes (\emph{estratégias}) para alinhar as sequências. Alguns são mais rápidos e menos
precisos, ideais para um grande número de sequências ou para fragmentos codificantes, sem
estrutura secundária; outros são mais precisos e mais lentos, ideais, por exemplo, para
fragmentos com estrutura secundária, como o 16S. Por serem mais lentos, estes últimos não podem
ser utilizados em grandes conjuntos de sequências. As duas estratégias que usaremos são:

\begin{center}
\begin{tabular}{@{}lll@{}}
\toprule
Estratégia & Opções no comando & Uso nesta prática \\
\midrule
FFT-NS-2 (rápida, padrão) & \verb|--retree 2| & COI \\
E-INS-i (precisa, lenta) & \verb|--genafpair --maxiterate 1000| & 16S \\
\bottomrule
\end{tabular}
\end{center}

A FFT-NS-2 parte de uma matriz de distâncias confiável, mas não é tão lenta quanto outras
estratégias do MAFFT. A E-INS-i lida bem com fragmentos como o 16S, em que regiões conservadas
se alternam com regiões cheias de inserções e deleções; ela faz iterações, refinando o
alinhamento até chegar a um resultado satisfatório, e \verb|--maxiterate 1000| aumenta o
número de iterações possíveis para 1000. Em todos os comandos usaremos também
\verb|--preservecase|, para que as bases sejam mantidas em letra maiúscula. Para saber mais
sobre as estratégias, veja
\url{https://mafft.cbrc.jp/alignment/software/algorithms/algorithms.html}.

\subsection{Instalação}

\paragraph{Windows (Ubuntu no WSL).} No terminal do Ubuntu, digite os comandos abaixo; o
\verb|sudo| pede a senha do Ubuntu (que não aparece enquanto é digitada) e o \verb|apt|
pergunta se deve continuar (responda \verb|Y|):
\begin{Verbatim}
$ sudo apt update
$ sudo apt install mafft
\end{Verbatim}

\paragraph{macOS.} Na página do MAFFT (\url{https://mafft.cbrc.jp/alignment/software/}),
entre em \ui{Mac OS X}, baixe o pacote padrão (\arq{.pkg}), abra-o e siga as instruções do
instalador. Quem usa o Homebrew pode, em vez disso, digitar \verb|brew install mafft|.

\paragraph{Conferindo.} Feche e abra o terminal e digite:
\begin{Verbatim}
$ mafft --version
\end{Verbatim}
Deve aparecer o número da versão (por exemplo, \verb|v7.505|, seguido da data). Lembre-se de que
o símbolo \verb|$| representa o \emph{prompt} e não deve ser digitado; no seu terminal ele
pode ser outro, como \verb|%|.

\subsection{Indo até a pasta da prática}

O MAFFT precisa encontrar os arquivos, então primeiro vamos até a pasta
\arq{pratica_alinhamento} com o \verb|cd| (Prática~2).
\begin{itemize}
\item \textbf{Windows:} os arquivos do Windows ficam, no Ubuntu, dentro de \arq{/mnt/c}.
Se a pasta estiver em \ui{Downloads}, por exemplo, o comando é o abaixo, trocando
\verb|Nome| pelo seu usuário do Windows (a tecla \tecla{Tab} completa os nomes):
\begin{Verbatim}
$ cd /mnt/c/Users/Nome/Downloads/pratica_alinhamento
\end{Verbatim}
Outra opção: abra a pasta no
Explorador de Arquivos, clique na barra de endereço, digite \verb|wsl| e tecle \tecla{Enter};
o Ubuntu abre já dentro da pasta.
\item \textbf{macOS:} digite \verb|cd| seguido de um espaço, arraste a pasta do Finder para
a janela do Terminal e tecle \tecla{Enter}.
\end{itemize}
Confira se você está no lugar certo: \verb|pwd| mostra o caminho da pasta atual, e a lista
mostrada por \verb|ls| deve incluir as pastas \arq{R} e \arq{sequencias}.
\begin{Verbatim}
$ pwd
$ ls
\end{Verbatim}
Agora crie a pasta onde ficarão os alinhamentos:
\begin{Verbatim}
$ mkdir alinhamentos
\end{Verbatim}

\subsection{Alinhando um marcador}

Vamos começar pelo COI, com a estratégia FFT-NS-2:
\begin{Verbatim}
$ mafft --preservecase --retree 2 \
      sequencias/COI_ids.fasta > alinhamentos/COI_alinhado.fasta
\end{Verbatim}
A barra invertida \verb|\| no fim da primeira linha indica que o comando continua na linha
de baixo: tecle \tecla{Enter} logo depois dela e continue digitando (você também pode
escrever tudo em uma única linha, sem a barra). Por partes:
\begin{itemize}
\item \verb|mafft --preservecase --retree 2|: o programa e as opções (Seção~\ref{sec:mafft});
\item \arq{sequencias/COI_ids.fasta}: o arquivo de entrada, com as sequências desalinhadas;
\item \verb|> alinhamentos/COI_alinhado.fasta|: o sinal \verb|>| redireciona a saída do
MAFFT, isto é, o alinhamento, para um arquivo, como na Prática~2. \textbf{Não use o mesmo nome
do arquivo de entrada}, ou ele será apagado antes de ser lido!
\end{itemize}
Enquanto trabalha, o MAFFT mostra mensagens de progresso na tela; elas não vão para o arquivo.
Quando o \emph{prompt} voltar, confira com \verb|ls alinhamentos| se o arquivo foi criado.

\subsection{Alinhando os dois marcadores de uma vez}\label{sec:laco}

Para o 16S, bastaria repetir o comando trocando o nome dos arquivos e a estratégia. Com muitos
marcadores, porém, repetir comandos dá trabalho e abre espaço para erros. Um \textbf{laço}
(\emph{loop}) executa os mesmos comandos para cada item de uma lista. Digite as linhas abaixo,
teclando \tecla{Enter} ao fim de cada uma:
\begin{Verbatim}
for marcador in 16S COI
do
  if [ "$marcador" = "16S" ]
  then
    mafft --preservecase --genafpair --maxiterate 1000 \
      sequencias/${marcador}_ids.fasta > alinhamentos/${marcador}_alinhado.fasta
  else
    mafft --preservecase --retree 2 \
      sequencias/${marcador}_ids.fasta > alinhamentos/${marcador}_alinhado.fasta
  fi
done
\end{Verbatim}
Depois da primeira linha, o terminal troca o \emph{prompt} por um sinal de continuação
(\verb|>| no bash do Ubuntu, \verb|for>| no zsh do macOS), indicando que espera o resto do
laço. Os comandos só são executados depois que você digita \verb|done|. Os espaços no início
das linhas são opcionais: servem apenas para facilitar a leitura.

Como o laço funciona:
\begin{itemize}
\item \verb|for marcador in 16S COI| cria a variável \verb|marcador|, que vale \verb|16S|
na primeira volta e \verb|COI| na segunda. Tudo o que está entre \verb|do| e \verb|done| é
repetido a cada volta.
\item \verb|${marcador}| é substituído pelo valor da variável: na primeira volta,
\arq{sequencias/${marcador}_ids.fasta} vira \arq{sequencias/16S_ids.fasta}. As chaves são
necessárias porque, sem elas, o terminal leria \verb|$marcador_ids| como o nome de outra
variável, \verb|marcador_ids|, que não existe.
\item \verb|if [ "$marcador" = "16S" ]| testa se o marcador é o 16S. Se for (\verb|then|),
usa a E-INS-i; se não (\verb|else|), usa a FFT-NS-2. O \verb|fi| (\verb|if| ao contrário)
fecha o teste. Os espaços dentro dos colchetes são obrigatórios.
\end{itemize}

O alinhamento do 16S leva mais tempo; espere o \emph{prompt} voltar e confira:
\begin{Verbatim}
$ ls alinhamentos
16S_alinhado.fasta  COI_alinhado.fasta
\end{Verbatim}
O \arq{COI_alinhado.fasta} foi refeito pelo laço, com o mesmo resultado do comando anterior.

Abra os dois alinhamentos no MEGA (\ui{Align > Edit/Build Alignment > Retrieve sequences from
a file}) ou em um editor de texto e compare-os: \textbf{qual a principal diferença entre o
alinhamento do 16S e o do COI? Qual dos dois demorou mais para ser concluído? Por quê?}
Para medir o tempo, repita um dos comandos com \verb|time| na frente (por exemplo,
\verb|time mafft --preservecase --retree 2 ...|).

\begin{resposta}
O COI (cerca de 712 colunas com o MAFFT v7.471) só tem \emph{gaps} no início e no fim,
causados pelas diferenças de comprimento das sequências: não há inserções nem deleções
internas, pois em um gene codificante elas mudariam a fase de leitura (a não ser que fossem
múltiplas de três). O 16S (cerca de 2475 colunas) tem muitos \emph{gaps} internos, porque é
um gene de RNA ribossomal, com regiões (as alças da estrutura secundária) que acumulam
inserções e deleções. O 16S demorou mais (cerca de 6~s contra 0,2~s, em um computador de dois
núcleos): usa a E-INS-i, com até 1000 iterações de refinamento, e tem mais sequências (22
contra 14), algumas delas longas (até cerca de 2400~pb).
\end{resposta}

% =====================================================================
\section{Concatenando os alinhamentos no R}\label{sec:concatenar}
% =====================================================================

Nas próximas práticas, os dois marcadores serão analisados juntos, em uma única matriz: o 16S
nas primeiras colunas e o COI em seguida. Essa união se chama \textbf{concatenação}, e deve
ser feita \textbf{pelo nome da espécie}: a linha de \emph{Hylodes nasus} da matriz final deve
juntar o 16S e o COI de \emph{H.~nasus}, embora os dois fragmentos tenham números de acesso
diferentes. Quando uma espécie tem apenas um dos marcadores, o outro é preenchido com
\verb|?|, o símbolo de \textbf{dado faltante}.

Abra o RStudio, crie um projeto na pasta \arq{pratica_alinhamento} (\ui{File > New Project >
Existing Directory}) e abra o script \arq{R/04_concatenar.R}. Execute-o linha a linha
(\tecla{Ctrl+Enter}, ou \tecla{Cmd+Enter} no macOS), acompanhando os resultados no console.
O script usa o pacote ape \cite{paradis2019}, instalado na Prática~3. Os principais passos
são descritos abaixo.

\paragraph{1. Ler os alinhamentos.} \verb|read.FASTA()| lê o arquivo e \verb|as.matrix()|
transforma o alinhamento em uma matriz, com uma linha por sequência e uma coluna por posição;
\verb|dim()| mostra o número de linhas e de colunas:
\begin{Verbatim}
s16 <- as.matrix(read.FASTA("alinhamentos/16S_alinhado.fasta"))
coi <- as.matrix(read.FASTA("alinhamentos/COI_alinhado.fasta"))
dim(s16)
dim(coi)
\end{Verbatim}

\paragraph{2. Separar a espécie do número de acesso.} Duas funções com expressões regulares
(Prática~1) separam as partes do nome:
\begin{Verbatim}
especie <- function(x) sub("_[A-Z]+_?[0-9]+$", "", x)
acesso  <- function(x) sub("^.*_([A-Z]+_?[0-9]+)$", "\\1", x)
\end{Verbatim}
Assim, \verb|Crossodactylus_schmidti_AY843579| vira \verb|Crossodactylus_schmidti| e
\verb|AY843579|. O script monta uma tabela com o acesso usado para cada espécie em cada
marcador e, em seguida, renomeia as linhas das duas matrizes apenas com o nome da espécie.

\paragraph{3. Completar com dados faltantes.} \verb|union()| reúne os nomes das duas
matrizes, e \verb|setdiff()| mostra quais espécies faltam em cada uma:
\begin{Verbatim}
nomes <- sort(union(rownames(s16), rownames(coi)))
setdiff(nomes, rownames(s16))   # espécies sem 16S
setdiff(nomes, rownames(coi))   # espécies sem COI
\end{Verbatim}
A função \verb|completar()| acrescenta a cada matriz uma linha de \verb|?| para cada espécie
que falta e põe as linhas na mesma ordem (a de \verb|nomes|).

\paragraph{4. Concatenar.} Com as duas matrizes com as mesmas linhas, na mesma ordem,
\verb|cbind()| junta as colunas:
\begin{Verbatim}
concatenada <- cbind(s16_completo, coi_completo)
dim(concatenada)
\end{Verbatim}
O número de colunas da matriz concatenada deve ser a soma das colunas do 16S e do COI.

\paragraph{5. Gravar.} O script cria a pasta \arq{matriz} com:
\begin{itemize}
\item \arq{hylodidae_concatenado.fasta}, \arq{.phy} e \arq{.nex}: a mesma matriz nos formatos
FASTA, PHYLIP e NEXUS, pedidos por diferentes programas de análise filogenética;
\item \arq{particoes.txt}: as colunas ocupadas por cada fragmento, que serão usadas para
analisar cada marcador com seu próprio modelo de evolução:
\begin{Verbatim}
DNA, 16S = 1-2475
DNA, COI = 2476-3187
\end{Verbatim}
\item \arq{tabela_acessos.csv}: os números de acesso usados para cada espécie, a tabela que
acompanha a publicação de um trabalho.
\end{itemize}

Os números de colunas citados acima são os obtidos com o MAFFT v7.471; com outra versão,
eles podem variar um pouco.

% =====================================================================
\section*{Para refletir}
\addcontentsline{toc}{section}{Para refletir}
% =====================================================================

\begin{enumerate}
\item No MEGA, por que foi preciso trocar o código genético para \ui{Vertebrate
Mitochondrial}? O que aconteceria com o código \ui{Standard}?
\item O que seria preciso mudar no laço da Seção~\ref{sec:laco} para alinhar também um
terceiro marcador, por exemplo o citocromo~\emph{b}, em um arquivo \arq{cytb_ids.fasta}, com
a estratégia FFT-NS-2?
\item Quantas espécies há na matriz concatenada? Quais delas não têm o 16S e quais não têm o
COI? Quantas colunas tem a matriz?
\item Por que a concatenação é feita pelo nome da espécie, e não pela posição das sequências
nos arquivos ou pelo número de acesso?
\end{enumerate}

\begin{resposta}
1. No código mitocondrial de vertebrados, \texttt{TGA} codifica triptofano (e não é códon de
parada), \texttt{ATA} codifica metionina, e \texttt{AGA} e \texttt{AGG} são códons de parada.
Com o código \ui{Standard}, os \texttt{TGA} apareceriam como \texttt{*}, sugerindo,
erradamente, que as sequências têm \emph{stop codons}.
2. Basta acrescentar o marcador à lista: \verb|for marcador in 16S COI cytb|. Como o teste
\verb|if| só separa o 16S, o \verb|cytb| cai no \verb|else| e é alinhado com a FFT-NS-2.
3. 24 espécies (as 22 do 16S mais \emph{Hylodes perere} e \emph{Phantasmarana apuana}, que só
têm COI). Sem 16S: \emph{H.~perere} e \emph{P.~apuana}. Sem COI: \emph{Crossodactylus aeneus},
\emph{C.~bokermanni}, \emph{C.~werneri}, \emph{Hylodes caete}, \emph{H.~charadranaetes},
\emph{H.~fredi}, \emph{H.~lateristrigatus}, \emph{H.~meridionalis}, \emph{H.~uai} e
\emph{Phantasmarana jordanensis}. A matriz tem 3187 colunas (2475 + 712) com o MAFFT v7.471.
4. As sequências de cada marcador vêm de registros diferentes, com números de acesso
diferentes (por exemplo, \emph{H.~nasus}: \texttt{KJ961577} no 16S e \texttt{KJ961557} no
COI), e os arquivos têm números diferentes de sequências, em ordens diferentes. O que liga os
dois fragmentos é o organismo; juntá-los pela posição misturaria espécies diferentes na mesma
linha.
\end{resposta}

% =====================================================================
\clearpage
\appendix
\section{Problemas comuns}\label{ap:problemas}
% =====================================================================

\begin{description}[style=nextline,leftmargin=1.5em]
\item[\texttt{mafft: command not found}] O MAFFT não foi instalado, ou o terminal foi aberto
antes da instalação: feche-o e abra-o de novo. No macOS, com o pacote \arq{.pkg}, o programa
fica em \arq{/usr/local/bin/mafft}.
\item[\texttt{No such file or directory}] Você não está na pasta \arq{pratica_alinhamento},
ou o nome do arquivo foi digitado errado. Confira com \verb|pwd| e \verb|ls|.
\item[Arquivo alinhado vazio (0 bytes)] O MAFFT não conseguiu ler a entrada (por exemplo,
porque o nome do arquivo de entrada está errado); leia a mensagem de erro na tela. Se o erro
for sobre a pasta \arq{alinhamentos}, ela ainda não existe: rode \verb|mkdir alinhamentos|.
\item[O terminal fica mostrando \texttt{>} e não executa nada] O laço ou o \verb|if| não foi
fechado: falta digitar \verb|fi| ou \verb|done|. Se errar no meio, tecle \tecla{Ctrl+C} e
comece de novo.
\item[\texttt{cannot open file} no R] A pasta de trabalho não é \arq{pratica_alinhamento}, ou
os alinhamentos ainda não foram feitos. Confira com \verb|getwd()| e \verb|list.files()|.
\item[\texttt{there is no package called \textquotesingle ape\textquotesingle}] Instale o pacote como na Prática~3
(\verb|install.packages("ape")|).
\end{description}

% =====================================================================
\section{Alterações em relação à versão anterior}\label{ap:mudancas}
% =====================================================================

\begin{enumerate}
\item \textbf{Sequências.} No lugar de \arq{hylodidae_16S.fas} e \arq{hylodidae_coi.fas},
são usadas as matrizes baixadas por número de acesso na Prática~3, com nomes no formato
\verb|Genero_especie_acesso| e a taxonomia atual (\emph{Phantasmarana}).
\item \textbf{MEGA.} Atualização do MEGA7 para o MEGA 12 \cite{kumar2024}; apenas o COI é
alinhado no MEGA. As capturas de tela foram substituídas por descrições em texto.
\item \textbf{MAFFT no terminal.} O menu interativo do \arq{mafft.bat} (Windows) foi
substituído por comandos digitados no terminal do Ubuntu (WSL) ou do macOS, com instruções de
instalação, e os dois marcadores passaram a ser alinhados com um laço.
\item \textbf{Opções do menu.} O texto anterior mandava escolher a ``opção dois'' para a
FFT-NS-2, que é a opção~3 do menu, e indicava a opção~6 (E-INS-i) para o 16S, mas o exemplo
mostrava a opção~3. As estratégias agora são indicadas pelas próprias opções de linha de
comando (\verb|--retree 2| e \verb|--genafpair --maxiterate 1000|).
\item \textbf{Ordem das sequências.} A saída ordenada por similaridade (\ui{Sorted},
\verb|--reorder|) foi substituída pela ordem de entrada, o padrão do MAFFT.
\item \textbf{Servidor do MAFFT.} A seção sobre o alinhamento no servidor \emph{on-line} foi
retirada.
\item \textbf{Concatenação.} Incluída a concatenação dos dois alinhamentos no R, pelo nome da
espécie, com dados faltantes (\verb|?|), arquivo de partições e tabela de números de acesso.
\item \textbf{Revisão do texto.} Correção de erros de digitação (``múltilos'', ``Edgard'',
``co relação'') e inclusão da seção \emph{Para refletir}.
\end{enumerate}

\begin{thebibliography}{9}
\bibitem{edgar2004} Edgar, R.~C. 2004. MUSCLE: multiple sequence alignment with high accuracy
and high throughput. \emph{Nucleic Acids Research} 32(5): 1792--1797.
\bibitem{katoh2013} Katoh, K. \& Standley, D.~M. 2013. MAFFT multiple sequence alignment
software version 7: improvements in performance and usability. \emph{Molecular Biology and
Evolution} 30(4): 772--780.
\bibitem{kumar2024} Kumar, S., Stecher, G., Suleski, M., Sanderford, M., Sharma, S. \&
Tamura, K. 2024. MEGA12: Molecular Evolutionary Genetic Analysis version 12 for adaptive and
green computing. \emph{Molecular Biology and Evolution} 41(12): msae263.
\bibitem{larkin2007} Larkin, M.~A., Blackshields, G., Brown, N.~P., Chenna, R., McGettigan,
P.~A., McWilliam, H., Valentin, F., Wallace, I.~M., Wilm, A., Lopez, R., Thompson, J.~D.,
Gibson, T.~J. \& Higgins, D.~G. 2007. Clustal W and Clustal X version 2.0.
\emph{Bioinformatics} 23(21): 2947--2948.
\bibitem{notredame2000} Notredame, C., Higgins, D.~G. \& Heringa, J. 2000. T-Coffee: a novel
method for fast and accurate multiple sequence alignment. \emph{Journal of Molecular Biology}
302(1): 205--217.
\bibitem{paradis2019} Paradis, E. \& Schliep, K. 2019. ape 5.0: an environment for modern
phylogenetics and evolutionary analyses in R. \emph{Bioinformatics} 35: 526--528.
\end{thebibliography}

\end{document}
